\bookchapter{الفصل 4: تثبيت التطبيقات وإدارتها — “حقيبة المهارات” لهاتفك}

\chapterintro{تجعل التطبيقات الهاتف قادراً على كل شيء، لكن يجب إدارتها جيداً. تماماً كما ترتيب مكان عمل لموظف جديد وتمنحه الصلاحيات، يحتاج تثبيت التطبيقات إلى عملية منظمة.}

\booksection{أولاً: المتاجر الرسمية: “سوبر ماركت التطبيقات” الأكثر أماناً}

\begin{itemize}
  \item \textbf{أندرويد}: Google Play (في الخارج) أو متاجر العلامات المدمجة (مثل Huawei AppGallery وXiaomi App Store).
  \item \textbf{iOS}: App Store فقط.
  \item \textbf{HarmonyOS}: Huawei AppGallery.
\end{itemize}

الأهمية الأمنية الأساسية: مراجعة، ضمان التحديثات، أمان الدفع. دليل العمل:\textbf{يفضّل التنزيل من المتاجر الرسمية فقط!}

\booksection{ثانياً: الوضع النقي/وضع الحماية: تفعيل “جدار الحماية” بهاتفك}

تدمج العلامات الصينية الرئيسية (Huawei وXiaomi وOPPO وvivo وHonor) جميعها\textbf{الوضع النقي}أو\textbf{وضع الحماية}، وهو مفتاح أمان على مستوى النظام.\textbf{يُوصى بشدة بتفعيله لجميع المستخدمين}.

\booksubsection{ما هو بالضبط؟}

الوظيفة الأساسية للوضع النقي (المعروف أيضاً بـ“وضع الحماية”) هي\textbf{منع تثبيت التطبيقات من مصادر خارجية}: عند تفعيله، يسمح النظام\textbf{بتثبيت التطبيقات من المتاجر الرسمية فقط}، وتُعترَض حزم التثبيت المنزّلة من المتصفحات والرموز البرمجية وبرامج الدردشة وغيرها من المصادر من المصدر، مما يمنع البرمجيات الخبيثة من البداية.

\booksubsection{الفئة المستهدفة}

\begin{itemize}
  \item \textbf{المستخدمون كبار السن}: الذين لا يفرّقون بين إعلانات "الرسمية" و"المزيفة"، الوضع النقي هو التعويذة الأكثر أماناً.
  \item \textbf{المستخدمون غير الحساسين لأذونات التطبيقات}: الذين لا يحتاجون إلى تثبيت تطبيقات متخصصة أو نسخ تجريبية أو ألعاب بخدمات أجنبية.
  \item \textbf{جميع المستخدمين الراغبين في رفع أمان أجهزتهم}: حتى للمستخدمين المتقدمين، التفعيل أفضل من عدمه — عند الحاجة المؤقتة، يمكن إيقافه مؤقتاً.
\end{itemize}

\booksubsection{مسار التفعيل (بأخذ Huawei/Xiaomi كمثال)}

\begin{itemize}
  \item \textbf{Huawei/HarmonyOS}: [الإعدادات] > [النظام والتحديث] >\textbf{[الوضع النقي]}، اختر "تفعيل".
  \item \textbf{Xiaomi}: [الإعدادات] > [حماية الخصوصية] >\textbf{[حماية الخصوصية]}أو [الإعدادات] > [إعدادات التطبيقات] > [تثبيت التطبيقات]، أوقف "تثبيت التطبيقات غير المعروفة".
  \item \textbf{OPPO/vivo}: [الإعدادات] > [الأمان] >\textbf{[تثبيت التطبيقات من مصادر خارجية]}، أوقف التفويض لكل تطبيق.
\end{itemize}

\booksubsection{آلية الفحص الأمني في متجر التطبيقات}

متجر التطبيقات ليس مجرد “ديكور”، بل خلفه خط إنتاج أمني متكامل:

\begin{itemize}
  \item \textbf{فحص الفيروسات}: يخضع كل تطبيق قبل طرحه لفحص مشترك من عدة محركات مكافحة فيروسات، للتعرف على بصمات أحصنة طروادة والفيروسات المعروفة.
  \item \textbf{مراجعة الأذونات}: يراجع المتجر ما إذا كانت الأذونات التي يطلبها التطبيق "معقولة" — تطبيق مصباح يدوي يطلب إذن الوصول إلى جهات الاتصال؟ يُرفض أو يُسحب فوراً.
  \item \textbf{توثيق المطوّر}: يتطلب الطرح توثيق مؤسسي أو بطاقة هوية شخصية أو بطاقة ائتمان. عند حدوث مشكلة يمكن تحديد المسؤول، وليس "تطبيقاً مظلماً".
  \item \textbf{المراقبة المستمرة}: تُعاد مراجعة التطبيقات المطروحة بشكل مستمر، وتُسحب التطبيقات التي تُكتشف سلوكياتها الخبيثة.
\end{itemize}

\booksection{ثالثاً: المصادر الخارجية و“الطفيليات الإلكترونية”: مخاطر عالية جداً!}

تنزيل التطبيقات من صفحات الويب أو روابط غير معروفة قد يحمل فيروسات، ويسرق الخصوصية، ويسحب أموال الدفع. هنا نشير بشكل خاص إلى مفهوم جديد:\textbf{“الطفيليات الإلكترونية”}.

\textbf{ما هي “الطفيليات الإلكترونية”؟}\par

تشير إلى\textbf{البرمجيات “المكسورة” المزروعة بأحصنة طروادة}— تظن أنك “وفّرت المال”، لكنك في الحقيقة “وقعت في الفخ”. الصور الشائعة:

\begin{itemize}
  \item برمجيات مدفوعة بـ"نسخة مكسورة" أو "نسخة مجانية" أو "نسخة بدون إعلانات"؛
  \item أدوات رمادية مثل"تضخيم الإعجابات" و"الاستيلاء التلقائي على المغلفات الحمراء" و"غش الألعاب"؛
  \item تطبيقات"النسخة الداخلية" المزعومة التي يوصي بها"الإعلانات الصغيرة" في المجموعات.
\end{itemize}

هذه البرامج قد تعمل ظاهرياً، لكنها خلف الكواليس قد\textbf{تسرق جهات الاتصال، وتقرأ الرسائل النصية، وتتنصت على الميكروفون، وتعدين العملات الرقمية في الخلفية}. الوضع النقي هو\textbf{الخط الدفاعي الأول}ضدها — يمنعك من تثبيتها أصلاً.

\booksubsection{توصيات الخط الدفاعي الثاني:}

\begin{itemize}
  \item \textbf{الفحص الشامل الدوري للفيروسات}: استخدم مركز الأمان المدمج في الهاتف (مثل Huawei Phone Manager وXiaomi Security Center) للفحص مرة كل شهر.
  \item \textbf{الحذر في منح الأذونات الحساسة}: جهات الاتصال، والرسائل النصية، والموقع، والكاميرا، والميكروفون — هذه الأذونات الخمس"عالية الخطورة" يجب مراجعتها بصرامة.
  \item \textbf{مراقبة استهلاك البطارية والبيانات غير الطبيعي}: سخونة مفاجئة في الهاتف، وانخفاض سريع في البطارية، واستهلاك البيانات في منتصف الليل — غالباً ما يعني وجود"طفيلي" مخفي في الخلفية.
\end{itemize}

إذا اضطررت لتثبيت تطبيق من مصدر خارجي:

\begin{enumerate}
  \item تأكد من أن المصدر موثوق تماماً.
  \item راجع متطلبات الأذونات بعناية قبل التثبيت.
  \item افحصه بتطبيق أمان بعد التثبيت.
\end{enumerate}

التوصية النهائية: في إعدادات النظام، أوقف إذن“تثبيت التطبيقات غير المعروفة”.

\booksection{رابعاً: التثبيت والتحديث وإلغاء التثبيت}

\begin{itemize}
  \item \textbf{التثبيت}: اضغط"الحصول" أو "التثبيت" في المتجر، وعند الفتح الأول سيُطلب منك الإذن.
  \item \textbf{التحديث}:\textbf{يجب التحديث!}أساساً لإصلاح الثغرات الأمنية. يُستحسن تفعيل التحديث التلقائي.
  \item \textbf{إلغاء التثبيت}: اضغط مطولاً على الأيقونة واسحبها إلى منطقة إلغاء التثبيت. للتنظيف الأكثر شمولاً، اذهب إلى [الإعدادات] > [إدارة التطبيقات] وابحث عنه وألغِ تثبيته.
\end{itemize}

\booksection{خامساً: إدارة الأذونات: أنت“المدير” للتطبيق}

طلب التطبيقات للأذونات (مثل جهات الاتصال والموقع) أحياناً يكون معقولاً (الخرائط تحتاج الموقع)، وأحياناً غير معقول (المصباح اليدوي يطلب جهات الاتصال).

\begin{itemize}
  \item \textbf{خيارات الإذن}:"السماح دائماً"،"السماح أثناء الاستخدام فقط"،"ممنوع"،"السؤال في كل مرة".
  \item \textbf{مبادئ الإدارة}:
  \begin{enumerate}
    \item التقليل: لا تمنح ما ليس ضرورياً.
    \item التدقيق اللاحق: راجع بانتظام إعدادات الأذونات لترى من"يتجاوز صلاحياته".
    \item السياق: مثل تطبيق استدعاء السيارة، امنح إذن الموقع فقط عند الاستخدام.
  \end{enumerate}
\end{itemize}

\begin{tcbwarn}{أذونات تحتاج إلى مراقبة دقيقة}
\begin{itemize}
  \item جهات الاتصال/الرسائل النصية: قد تُستخدم لإرسال رسائل عشوائية.
  \item الموقع الفوري: قد يُستخدم لتتبع تحركاتك.
  \item الكاميرا/الميكروفون: قد يُستخدم للتصوير والتسجيل الخفي.
  \item الوصول إلى الملفات: قد يقرأ ملفاتك الخاصة.
\end{itemize}

\end{tcbwarn}

\booksection{سادساً: مكونات حزمة تثبيت البرنامج}

هل تساءلت يوماً عمّا تحتويه حزمة التثبيت التي تنزّلها من متجر التطبيقات؟ فهم مكونات حزمة التثبيت يساعدك على فهم كيفية عمل التطبيقات بشكل أفضل، ومعرفة من أين تبدأ عند مواجهة مشكلات.

\booksubsection{بنية ملف APK في Android}

APK (Android Package Kit) هو في جوهره حزمة مضغوطة بصيغة ZIP، ويمكنك بعد فك ضغطه رؤية المجلدات والملفات الأساسية التالية:

\begin{itemize}
  \item \textbf{AndroidManifest.xml}: "بطاقة هوية" التطبيق و"دليله". يسجل اسم حزمة التطبيق ورقم الإصدار والأذونات المطلوبة وإصدارات النظام المدعومة ونقطة الدخول Activity وغيرها من المعلومات الأساسية. يقرأ النظام هذا الملف أولاً عند تثبيت التطبيق.
  \item \textbf{classes.dex}: الكود الأساسي للتطبيق. يتحول كود المصدر Java/Kotlin بعد التجميع إلى Dalvik bytecode ويُخزَّن في هذا الملف. قد يحتوي التطبيق على عدة ملفات dex (مثل classes2.dex وclasses3.dex).
  \item \textbf{resources.arsc}: جدول فهرس الموارد المجمّع. يسجل معرّفات جميع الموارد والعلاقات المقابلة لها.
  \item \textbf{مجلد res/}: مجلد الموارد. يحتوي على الصور وملفات التخطيط والسلاسل وتعريفات الألوان وغيرها. تُوضع صور الدقات المختلفة في مجلدات drawable مختلفة (مثل drawable-hdpi وdrawable-xhdpi).
  \item \textbf{مجلد assets/}: مجلد الموارد الخام. يخزن الموارد غير المجمّعة، مثل ملفات الويب والخطوط والملفات الصوتية. يمكن للتطبيق قراءتها عبر AssetManager عند التشغيل.
  \item \textbf{مجلد lib/}: مجلد المكتبات المحلية. يخزن ملفات .so المجمّعة (ناتج تجميع كود C/C++)، مصنّفة حسب بنية المعالج (مثل arm64-v8a وarmeabi-v7a). عادة ما تحتوي الألعاب والتطبيقات التي تتطلب حسابات عالية الأداء على هذا المجلد.
  \item \textbf{مجلد META-INF/}: مجلد معلومات التوقيع. يحتوي على ثلاثة ملفات: MANIFEST.MF وCERT.SF وCERT.RSA، للتحقق من سلامة APK وتوقيعه.
\end{itemize}

\booksubsection{بنية ملف IPA في iOS}

IPA (iOS App Store Package) هو أيضاً حزمة مضغوطة، لكن بنيته تختلف قليلاً:

\begin{itemize}
  \item \textbf{مجلد Payload/}: حيث يوجد المحتوى الأساسي. يحتوي بداخله على مجلد .app يضم جميع محتويات التطبيق.
  \item \textbf{Info.plist}: ملف إعدادات التطبيق، مشابه لـ AndroidManifest.xml في Android. يسجل معرّف التطبيق والإصدار ومتطلبات الأذونات والأجهزة المدعومة وغيرها.
  \item \textbf{الملف التنفيذي}: البرنامج الرئيسي الذي يحمل اسم مجلد .app، وهو ملف ثنائي بصيغة Mach-O المجمّعة.
  \item \textbf{مجلد Frameworks/}: المكتبات الديناميكية والأطر التي يعتمد عليها التطبيق.
  \item \textbf{ملفات الموارد}: الصور والصوت والفيديو وغيرها من الموارد، عادة ما تُخزَّن داخل مجلد .app.
  \item \textbf{مجلد \_CodeSignature/}: معلومات توقيع الكود، للتحقق من سلامة التطبيق ومصدره.
\end{itemize}

\booksubsection{آلية التحقق من التوقيع}

لماذا التطبيقات المنزّلة من المتاجر الرسمية أكثر أماناً؟ التحقق من التوقيع هو المفتاح.

\begin{itemize}
  \item \textbf{توقيع Android}: يوقّع المطوّر ملف APK بمفتاح خاص، ويتحقق النظام عند التثبيت بالمفتاح العام. إذا تم التلاعب بـ APK (مثل زرع فيروس)، يصبح التوقيع غير صالح، ويرفض النظام التثبيت. كما يتطلب Google Play توقيع التطبيقات App Signing، حيث يحتفظ Google بمفاتيح التوقيع، مما يحمي أمان التطبيق بشكل إضافي.
  \item \textbf{توقيع iOS}: مراجعة Apple أكثر صرامة. يجب أن يحمل كل IPA توقيع شهادة صادرة من Apple، وإلا لا يمكن تثبيته. لهذا السبب لا يمكن تثبيت التطبيقات على iOS إلا من App Store أو عبر شهادات المؤسسات. التطبيقات الموقّعة بشهادات المؤسسات لها صلاحية محدودة، ويمكن لـ Apple إبطالها في أي وقت.
  \item \textbf{وظائف التوقيع}:
  \begin{enumerate}
    \item التحقق من المصدر: التأكد من أن التطبيق يأتي فعلاً من المطوّر المُدّعى.
    \item التحقق من السلامة: ضمان عدم التلاعب بالتطبيق أثناء النقل.
    \item منع الانتحال: لا يمكن للتطبيقات الخبيثة انتحال شخصية تطبيقات أخرى.
  \end{enumerate}
\end{itemize}

\booksubsection{حجم حزمة التثبيت والتحسين}

قد تلاحظ أن حجم التطبيق المعروض في متجر التطبيقات يختلف عن المساحة التي يشغلها بعد التثبيت. والسبب:

\begin{itemize}
  \item \textbf{حجم التنزيل}: حجم APK/IPA بعد الضغط.
  \item \textbf{حجم التثبيت}: الحجم بعد فك الضغط، وعادة ما يكون أكبر بكثير من حجم التنزيل.
  \item \textbf{الاستهلاك أثناء التشغيل}: يُنشئ التطبيق عند التشغيل ذاكرة تخزين مؤقت وملفات بيانات، فيشغل مساحة أكبر.
\end{itemize}

وسائل تحسين التطبيقات الحديثة:

\begin{itemize}
  \item \textbf{App Bundle (Android)}: تقنية التوزيع الديناميكي من Google Play. يحصل المستخدم عند التنزيل على الموارد المناسبة لجهازه فقط (مثل صور بدقة الشاشة المقابلة فقط)، مما يقلل حجم التنزيل بشكل كبير.
  \item \textbf{App Thinning (iOS)}: تقنية مشابهة، حيث تُقصّ Apple تلقائياً الموارد غير الضرورية وفقاً للجهاز.
  \item \textbf{التنزيل عند الطلب}: يمكن للتطبيقات الكبيرة مثل الألعاب تنزيل الحزمة الأساسية أولاً، ثم تنزيل حزم الموارد لاحقاً.
\end{itemize}

\begin{tcbinfo}{معلومة صغيرة: لماذا تكون حزم تثبيت بعض التطبيقات كبيرة جداً؟}
تطبيقات الألعاب عادة ما تكون الأكبر، لأنها تحتوي على كميات كبيرة من القوام عالي الدقة والنماذج ثلاثية الأبعاد والملفات الصوتية. تطبيقات التواصل الاجتماعي أيضاً كبيرة نسبياً، لأنها تدعم وحدات وظيفية متنوعة. تطبيقات الأدوات عادة ما تكون الأصغر. إذا لاحظت أن تطبيق أدوات كبير بشكل غير عادي، فاحذر من احتمال تضمين وظائف غير ضرورية.

\end{tcbinfo}

\begin{tcbpractice}{تطبيق عملي: مراجعة أذونات 3 تطبيقات}
اختر 3 تطبيقات شائعة الاستخدام (مثل التواصل والتسوق والأدوات)، وادخل صفحة إدارة أذونات التطبيقات في [الإعدادات]، وراجعها عنصراً عنصراً، وأوقف الأذونات الحساسة غير الضرورية أو غيّرها إلى“السماح أثناء الاستخدام فقط”.

\end{tcbpractice}
